% --- ОБЩИЕ НАСТРОЙКИ ТЕКСТА ---
\onehalfspacing
\setlength{\parindent}{1.25cm} 

% Борьба с дырками в строках и переносами
\sloppy 
\hyphenpenalty=1000 
\tolerance=2000

% --- ПРИНУДИТЕЛЬНОЕ СОДЕРЖАНИЕ (ФИКС "ОГЛАВЛЕНИЯ") ---
\usepackage{etoolbox}
\addto\captionsrussian{%
  \renewcommand{\contentsname}{СОДЕРЖАНИЕ}%
}

% Центрирование и шрифт заголовка Содержания
\BeforeTOCHead[toc]{%
  \let\raggedsection\centering% Принудительно по центру
  \setkomafont{chapter}{\rmfamily\bfseries\normalsize}% Жирный и 14pt
}

% После оглавления возвращаем заголовки глав влево (но с твоим отступом)
\AfterTOCHead[toc]{%
  \let\raggedsection\relax%
}

% --- ШРИФТЫ ПУНКТОВ (Обычный, не жирный) ---
\setkomafont{chapterentry}{\normalfont\normalsize}
\setkomafont{chapterentrypagenumber}{\normalfont\normalsize}
\KOMAoptions{toc=chapterentrydotfill}

% --- НАСТРОЙКА ЗАГОЛОВКОВ (14pt жирный + отступ 1.25) ---
\KOMAoptions{numbers=noendperiod}
\setkomafont{disposition}{\rmfamily\bfseries\normalsize}

% Тот самый рабочий сдвиг заголовков
\makeatletter
\renewcommand{\sectionlinesformat}[4]{%
    \hspace{1.25cm}#3#4%
}
\renewcommand{\chapterlinesformat}[3]{%
    \hspace{1.25cm}#2#3%
}
\makeatother

\RedeclareSectionCommand[beforeskip=12pt, afterskip=6pt]{chapter}
\RedeclareSectionCommand[beforeskip=12pt, afterskip=6pt]{section}
\RedeclareSectionCommand[beforeskip=12pt, afterskip=6pt]{subsection}

% --- СПИСКИ ---
\setlist{
	leftmargin=0pt,
	nosep,
	wide=\parindent,
	labelsep=0.5em,
	listparindent=\parindent
}
\setlist[itemize,1]{label=--} 
\setlist[enumerate,1]{label=\arabic*.}

% --- ПОДПИСИ И ТАБЛИЦЫ ---
\captionsetup{justification=centering, singlelinecheck=false}
\DeclareCaptionLabelFormat{simple}{#1~#2}
\captionsetup{labelformat=simple}
\DeclareCaptionLabelSeparator{emdash}{ --- }
\captionsetup[table]{
	labelsep=emdash,
	justification=raggedright,
	singlelinecheck=false,
	position=top,
	skip=6pt
}

% --- ЛИСТИНГИ ---
\lstset{
    basicstyle=\ttfamily\footnotesize,
    breaklines=true,
    texcl=true,
    frame=none
}